\chapter{How to Train}\label{ch:iv1:how-to-train}

This chapter says what quantitative interviews test, in what order they test it, and how to use these three books to get ready. Read it once before you start, and come back when you plan a week. It holds no questions.

What follows about firms comes from two kinds of source: what firms publish about their own process, which is mostly logistics, and what candidates report in public, which is detailed but unaudited. So treat every number here as an indication rather than a specification. No firm is named, because a process varies by office, role and year, and a name would suggest a precision that nobody has.

\section{The stages}

Most processes have the same shape, with stages that run from the cheap and automatic to the expensive and personal.

\textbf{Application.} A CV, sometimes with a transcript or a short written motivation. Some firms invite an online test for every applicant whose CV passes a quick filter. One large market maker says that its process starts with an application and online assessments, followed by a behavioural interview with a recruiter.

\textbf{Online test.} One or more timed tests taken alone, before a person has looked at you closely. This is where reports say that most applications end, though no firm publishes the share. \Cref{sec:iv1:how-to-train:tests} describes what they contain.

\textbf{Recorded or recruiter screen.} A short call, or a video answered to a screen with no one on the other side. One large bank's student page describes a video interview of about 30 minutes. Candidates report a mix of motivation questions and some light mental arithmetic, which is partly a check on how you react.

\textbf{Technical rounds.} One to three calls of 30 to 60 minutes, by phone or video, with a person who works in the role. This is the core of this book and of the next two.

\textbf{Take-home work.} Not universal, but reported mainly at funds whose researchers work on data: a data exercise lasting anywhere from a few hours to two days, followed by a call that tests whether you understand your own work. For research analysts in fundamental investing, the reported equivalent is a stock pitch and a written case with a financial model, with about a week to complete it.

\textbf{Final round.} Several back-to-back interviews in one day, often called a superday. One large bank's student page says two to five for campus hires, depending on the division. Reports for engineers put it nearer four to six. Expect deeper versions of the technical rounds, some fit questions, and, at many trading firms, a market-making game.

\textbf{Decision.} Reported to take from a day or two to a week after the final round. A process for an intern or a new graduate is reported to take three to eight weeks overall.

\begin{dated}{2026-11}{Attempts are limited}\label{dat:iv1:how-to-train:retakes}
One large market maker's published recruitment page says that its online assessments can be taken only once every eight months, that there are no resets once an attempt has started, and that recent results carry over to a similar role. Check this for each firm. Do not take a real test for practice, and do not start one on a poor connection. Do your practice on the sheets in this book.
\end{dated}

\section{Online tests}\label{sec:iv1:how-to-train:tests}

Candidates report four ingredients, in varying mixes: mental arithmetic, probability and expected value, patterns and logic, and coding. Here are the formats they most often describe. All of them come from candidate reports gathered by preparation sites, and none is published by a firm.

\begin{dated}{2026-11}{Formats that candidates report}\label{dat:iv1:how-to-train:formats}
\textbf{Arithmetic against the clock.} The best-known format is 80 items in 8 minutes, with no calculator. Variants of 60 in 6 and 50 in 5 are also reported, some with four options per item and some with typed answers. Reports differ on the penalty: most say one mark lost for a wrong answer, and at least one says two. Candidates also report a pass mark, always unofficial, of about 55 on the 80-item version.

\textbf{Sequences.} Find the next term, with five options, in one reported case 26 items in 25 minutes, and in another 30 items in about 6 minutes.

\textbf{Probability and problem solving.} Reported examples are 15 to 17 questions in an hour, about 6 in 45 minutes, and 18 in 30 minutes. The questions take longer as the count falls.

\textbf{Cognitive games and questionnaires.} Short games that measure attention and working memory, sometimes with a personality questionnaire.

\textbf{Coding.} For developers and many researchers, typically two or three problems in 60 to 120 minutes on a hosted platform, marked by hidden test cases.
\end{dated}

The sheets in the speed-drills chapter (\cref{ch:iv1:speed-drills}) vary length, time and scoring so that you are not tuned to a single format. Before you take a real test, find out what you can about its scoring, and work out the guessing rule. A guess worth $+a$ when right and $-b$ when wrong pays only if you are right with probability above $b/(a + b)$. That single line decides how many items you should leave blank.

\section{What each role is asked}

The tags on each question show where it fits. This is the broad pattern behind them. Reports vary, and titles vary even more, so read the description of the job, not its title.

\textbf{Trader and quantitative trader.} Mental arithmetic, probability with no calculator, betting and market-making games. At the firms that train their own traders, candidates are expected to think quickly and aloud, and one large market maker says plainly that it will not test trading candidates on finance or economics. Bank and fund desks expect more knowledge of markets. Start with this book's speed and probability parts, then move on to the trading games in Book 19.

\textbf{Quantitative researcher.} Probability and statistics in depth, regression and machine learning, some coding, and a discussion of your own projects. PhD candidates are often asked to explain their research. This book's probability and statistics parts come first, then research cases in Book 19 and data handling in Book 20.

\textbf{Analyst.} Here the book means both data-minded analysts and fundamental analysts. A fundamental analyst at a fund is reported to face a stock pitch, long and short, in almost every interview, along with valuation and a model. Arithmetic and estimation from this book help, but the main work is in the valuation and pitch chapters of Book 19.

\textbf{Data scientist.} Statistics, regression, machine learning, experimental thinking and data handling, with coding in Python and SQL. Use the statistics part of this book and the data chapters of Book 20.

\textbf{Quantitative developer and software developer.} Algorithms and data structures, language depth, some probability and some system design. Book 20 is the main book, and this book's probability warm-ups are enough for most of the probability you will meet.

\textbf{Low-latency engineer.} Reports describe one or two live coding screens of about 60 minutes, followed by a day of four to six rounds on C++ internals, concurrency, cache behaviour and systems design, and sometimes a short quantitative round. Book 20 is the main book.

\textbf{Risk quant and bank quant.} Reports describe probability and statistics followed by pricing, stochastic calculus, C++ or Python, and, for validation roles, how to review someone else's model. Use this book's statistics and stochastic calculus chapters, and then the pricing and risk chapters of Book 19.

The type of firm shifts this too. Market makers lean towards speed and games, funds towards research depth and your own work, banks towards product knowledge, stochastic calculus and fit, and crypto firms towards the same mathematics plus market structure specific to that area. These are tendencies, and you will meet exceptions in every round.

\section{How to train}

Use the three books actively. The rest of this section is the book's method, and three studies support it.

\textbf{Answer before you read.} Commit to an answer, on paper or aloud, before you turn to the back. In a well-known experiment, students who tested themselves on a passage remembered more of it a week later than students who reread it, though the rereaders felt more confident. Reading an answer without having tried gives you that same false confidence.

\textbf{Space your practice.} A review of the research on spacing found that spreading practice over time beats cramming, and that the best gap grows with how long you need to remember the material. So take twenty minutes a day over six weeks rather than a long weekend. Return to a question you missed after a few days, then after a couple of weeks.

\textbf{Speak.} Most rounds are spoken. Say what you are doing, state your assumptions, give a rough answer early and refine it. An interviewer can help you only if they can hear where you are, and they are scoring how you reason as well as where you end up. If you train alone, record yourself answering a one-star question out loud, and listen once.

\textbf{Check every answer before you say it.} Test a small case, check the size and look at the units, and use a limiting case. The mental-technique chapter (\cref{ch:iv1:mental-technique}) gives quick checks for arithmetic. A wrong answer you catch yourself counts for you, and one that you announce with confidence does not.

\textbf{Practise under test conditions.} A meta-analysis of retesting found that cognitive test scores rise by about a quarter of a standard deviation on a second attempt, through practice and coaching. So practice does help on the format, and you want to have your practice effect before the real attempt. Use the drill sheets and the timed tests with a timer, in one sitting, and mark yourself honestly.

\section{A plan}

Stars give a target time (see the preface), and the tables in the preface count questions by chapter for each role. A sensible order for most candidates:

\begin{enumerate}
\item \textbf{Measure.} Take a drill sheet and the first timed test cold. This tells you what to work on, and it stops you from spending weeks on your strong side.
\item \textbf{Speed, a little every day.} Ten minutes on one drill sheet or one family a day, for the whole of your preparation. Speed needs regular practice, and a little each day costs almost nothing.
\item \textbf{Core topics.} Work through the chapters your role needs, in order within each chapter. Stop a chapter when the three-star questions are going well, and skip what your tags exclude. A question you solve in the star time counts as done. One you missed goes on a list to retry.
\item \textbf{Mixed tests.} From the middle of your preparation on, take one timed test a week. Questions sorted by topic give away the method. Timed tests do not, and they are the only honest measure of whether you have the method or only recognise it.
\item \textbf{The final stretch.} In the last two weeks, drop new material. Retry your list, take the remaining timed tests, and practise speaking through a mock interview from Book 19.
\end{enumerate}

If you have only a week, take a drill sheet a day, the first ten questions of each probability chapter and one timed test, and leave the rest. A book of this size is meant for a full preparation, and you do not have to finish it. It is better to do half of it with care than all of it with the answers open.
